\documentclass[a4paper, 12pt]{extarticle}
\usepackage[utf8]{inputenc}
\usepackage{soul}
\usepackage[ngerman]{babel}
\usepackage[T1]{fontenc}
\usepackage{graphicx}
\usepackage{scrpage2}
\usepackage{listings,xcolor,color}

%Seitenrand Einstellen
\usepackage[paper=a4paper,left=35mm,right=25mm,top=25mm,bottom=25 mm]{geometry}

% andere Schriftfont und Größe ändern
%\renewcommand{\familydefault}{\sfdefault}{\large}
%\usepackage{helvet}
%\large

% bibtex
\usepackage{cite}




\definecolor{mygreen}{rgb}{0,0.6,0}
\definecolor{mygray}{rgb}{0.5,0.5,0.5}
\definecolor{mymauve}{rgb}{0.58,0,0.82}



%opening
\title{Praxisbericht}
\author{Gennadi Heimann}
\date{20. September 2014}

\lstdefinestyle{bash}{
  language=sh,
  keywordstyle=\color{blue},
  commentstyle=\color{dkgreen},
  stringstyle=\color{mymauve},
  breaklines=true,
  numbers=left,
  basicstyle=\scriptsize\ttfamily,
  showspaces=false,
  frame=shadowbox,
  xleftmargin=0.5cm,
  tabsize=3
}

\lstdefinestyle{java} {
  language=java,
  basicstyle=\footnotesize,        % size of fonts used for the code
  breaklines=true,                 % automatic line breaking only at whitespace
  commentstyle=\color{mygreen},    % comment style
  escapeinside={\%*}{*)},          % if you want to add LaTeX within your code
  keywordstyle=\color{blue},       % keyword style
  stringstyle=\color{mymauve},     % string literal style
  numbers=left,
  tabsize=3,
  basicstyle=\small\ttfamily,
  showspaces=false,
  frame=shadowbox,
  xleftmargin=0.5cm,
}
\lstdefinestyle{sql} {
  language=SQL,
  basicstyle=\footnotesize,        % size of fonts used for the code
  breaklines=true,                 % automatic line breaking only at whitespace
  commentstyle=\color{mygreen},    % comment style
  escapeinside={\%*}{*)},          % if you want to add LaTeX within your code
  keywordstyle=\color{blue},       % keyword style
  stringstyle=\color{mymauve},     % string literal style
  numbers=left,
  tabsize=3,
  basicstyle=\small\ttfamily,
  showspaces=false,
  frame=shadowbox,
  xleftmargin=0.5cm
}

\begin{document}

\maketitle
\thispagestyle{empty}

\pagebreak

\pagestyle{scrheadings}
\chead{\includegraphics[scale=0.15]{img/doubleSlash_logo.jpg}}
\ohead{\today}
\ifoot{Gennadi Heimann}
\ofoot{\pagemark}
\setheadsepline{0.7pt}
\setfootsepline{0.7pt}

\tableofcontents

\pagebreak

\listoffigures

\pagebreak

\ihead{Headtitle}

\section{Beispiel}
Bla sagte Nobody \cite{felfernig1999intelligente}.
  
  \subsection{Eigene Produkte}
    
    \begin{itemize}
      \item \textbf{Beispiel}\footnote{url}
      \begin{itemize}
        \item
        \item 
      \end{itemize}
    \end{itemize}
  
    \subsubsection{Beispiel}
  

\pagebreak
\ihead{Headtitle}

\begin{figure}[htbp]
  \centering
  \includegraphics[scale=0.18]{./img/promt_verp.JPG}
  \caption {Caption}
\end{figure}
   
  \begin{lstlisting}[style=bash]
echo "#INFO# -> Sleeping for 10 seconds."
  \end{lstlisting}
  
  \emph{Kursive Beispiel}

\pagebreak
\ihead{Abbildungsverzeichnis}



\newpage
\ihead{Tagesberichte}

\begin{center}
  \begin{tabular}{|l|l|p{7cm}|}
    \hline
    Lf. Nr. & Wochentag, Datum & Tätigkeit\\
    \hline\hline
    1 & Mo. 3. März 2014 & Begrüßung. Fach- und GUI- Konzept lesen.\\
    \hline
    2 & Di. 4. März 2014 & Fach- und GUI- Konzept lesen.\\
    \hline
  \end{tabular}
\end{center}

\bibliography{thesis_bibtex}{}
%\bibliographystyle{alpha}
\bibliographystyle{apalike}

\end{document}
